\documentclass[10pt,a4paper]{amsart}
\usepackage[margin=19mm]{geometry}
\usepackage{amsmath,amssymb,mathtools}
\usepackage[scr=rsfs,cal=euler]{mathalfa}
\usepackage{xfrac}
\usepackage[unicode,hidelinks,bookmarks=false]{hyperref}
\newcommand{\ie}{i.e.\ }
\newcommand{\cf}{cf.\ }
\newcommand{\resp}{respectively\ }
\newcommand{\Subsection}{\subsection}
\providecommand{\nobreakdash}{\nobreak\hbox{-}}
\setlength{\emergencystretch}{2em}
\multlinegap0pt
\usepackage{amssymb}
\usepackage[normalem]{ulem}
\newtheorem{lemma}{Lemma}
\newcommand{\E}[1]{\mathbb{E}\lrs{#1}}
\newcommand{\R}{\mathbb{R}}
\newcommand{\lcs}{\ell_{cs}}
\newcommand{\lrp}[1]{\left( #1 \right)}
\newcommand{\lrs}[1]{\left[ #1 \right]}
\newcommand{\numberthis}{\addtocounter{equation}{1}\tag{\theequation}}
\newcommand{\xk}{x_k}
\title{Concentration of Stochastic Approximation Under Heavy-Tailed Markovian Noise}
\author{Shubhada Agrawal, Siva Theja Maguluri, Martin Zubeldia}
\date{Source: arXiv:2605.20999v1 (CC BY 4.0)}
\begin{document}
\maketitle

\noindent\emph{Source and licence.} Concentration of Stochastic Approximation Under Heavy-Tailed Markovian Noise. Version arXiv:2605.20999v1 (CC BY 4.0), distributed by arXiv under the Creative Commons Attribution 4.0 International licence, http://creativecommons.org/licenses/by/4.0/, as recorded in the arXiv licence metadata for this exact version and shown on the abstract page https://arxiv.org/abs/2605.20999v1. Full licence notice: \texttt{licenses/arxiv-2605.20999-CC-BY-4.0.txt}.\par\smallskip

\noindent\textit{Editorial scope.} Counted unit: the article's lemma lem:rec1beforeexp (source0.tex lines 1988--1993). The article's complete original proof of it is delivered with it (source0.tex lines 2196--2314, one \texttt{proof} environment of 119 lines, which the article places away from the statement). No mathematics is written, repaired or re-derived: the statement and the proof are the article's own lines, in the article's own notation, and no retained line is substituted or re-flowed.  Retained uncounted as the source-local context the proof needs: lines 1915--1919; lines 1920--1930; lines 1996--2002; lines 2540--2543.  Nothing was omitted from any retained span, so the package carries no omission record.  No retained line contains a citation, so the wrapper carries no bibliography and no bibliography entry.  No retained line contains a figure environment, an image-inclusion command or drawing code, so no image is delivered and the package declares no asset dependency.  Editorial formatting only, in addition: an AMS \texttt{amsart} preamble that restates the article's own declarations and macros used by the retained lines, namely the environments \texttt{lemma} on separate counters, together with the standard packages those lines need.  The retained environments print in this document as Lemma 1 in order of appearance, whereas the article prints the same items with the numbering of its own full document.  The complete native source of the article as distributed in the arXiv e-print for this exact version is bundled unchanged as \texttt{sources/201032/source0.tex}.
    \begin{align*}
        &M_{11} := \alpha_k \langle \nabla M(\xk - x^*), V_{x_k}(S_k) - \E{V_{x_k}(S_k) | S_{k-1}, x_{k-1}, Z_{k-1}} \rangle, \\
        & M_{12} := \alpha_k \langle \nabla M(\xk - x^*), \E{V_{x_k}(S_k) | S_{k-1}, x_{k-1}, Z_{k-1}} - \E{V_{x_k}(S_{k+1}) | S_{k}, x_{k}} \rangle, \numberthis \label{eq:T12}\\
        & M_{13} := \alpha_k \langle \nabla M(\xk - x^*), Z_k \rangle.
    \end{align*}

\begin{lemma}\label{lem:T11Properties}
    The terms $M_{11}$ and $M_{13}$, defined in Equation~\eqref{eq:T12}, satisfy 
    \begin{align*} 
    &\E{ M_{11} | S_{k-1}, x_{k-1}, Z_{k-1} }  = 0 \qquad \text{and} \qquad \E{M_{13} | S_{k-1}, x_{k-1}, Z_{k-1}} = 0.
    \end{align*}
    Moreover, they are bounded as below:
    \begin{align*} 
        &|M_{11}| \le \frac{2\alpha_k L_s L_1}{\lcs^2}\|\xk- x^*\|^2_c + \frac{2\alpha_k L_s L_2}{\lcs^2}\|\xk-x^*\|_c, \quad \text{and} \quad  |M_{13}| \le \alpha_k \frac{L_s B_2}{\lcs^2} \|x_k -x^*\|_c.
    \end{align*}
    Furthermore, $M_{11}$ and $M_{13}$ are conditionally independent, conditioned on $S_{k-1}, x_{k-1}, Z_{k-1}$. 
\end{lemma}

\begin{lemma}\label{lem:nonnegLyapunov} For $h \ge 8$ large enough so that for all $k\ge 1$
\[ 1 - \alpha_{k-1}\tfrac{uL_s(L_1+L_2)}{\lcs^2} \ge 0, \]
the following holds for all $k\ge 1$: 
    \begin{align*}
        M(x_k - x^*) + \alpha_{k-1} d_k + \tfrac{\alpha_{k-1}L_sL_2}{4} \ge 0.
    \end{align*}
\end{lemma}

\begin{lemma}\label{lem:ratioofalpha}
    For $\alpha > 0$, $h \ge 8$, $z\in (0,1]$, and for all $k\ge -1$,  $\alpha_k = \alpha/(k+h)^z$ or $\alpha_k = \alpha$, we have
    \[ \frac{\alpha_k}{\alpha_{k+1}}\le 2. \]
\end{lemma}

\begin{lemma}\label{lem:rec1beforeexp} The following bound holds.
    \begin{align*}
        M(x_{k+1} - x^*) + \alpha_k d_{k+1} &\le M(\xk - x^*) \lrp{1-\alpha_k \eta + \alpha^2_k C_1 } \\
        &\qquad+ M_{11} + M_{13} + \alpha_{k-1}d_k \lrp{1-\tfrac{\eta}{2}\alpha_k} + \alpha^2_k C_2.
    \end{align*}
\end{lemma}

\begin{proof} 
    Notice that $x_k$ is $\mathcal F_k = \sigma(x_{k-1},S_{k-1}, Z_{k-1})$ measurable since $x_k$ is a deterministic function of $x_{k-1}, S_{k-1}, $ and $Z_{k-1}$.
    
    In this proof, we define $\tilde{\gamma}_c := \eta/2$. Next, $\lambda_{k+1} \ge 0$, multiplying both sides of inequality in Lemma~\ref{lem:rec1beforeexp} by $\lambda_{k+1}$, exponentiating, and multiplying by, we get
    \begin{align*}
        e^{\lambda_{k+1}\lrp{M(x_{k+1} - x^*) + \alpha_k d_{k+1}}} 
        &\le e^{\lambda_{k+1} M(\xk - x^*) \lrp{1 - 2\alpha_k\tilde{\gamma}_c + \alpha^2_k C_1 }} \\
            &\qquad e^{\lambda_{k+1} (M_{11}+M_{13})} e^{\lambda_{k+1} \alpha_{k-1}d_k (1-\alpha_k \tilde{\gamma}_c) + \lambda_{k+1}\alpha^2_k C_2}.
    \end{align*}
    Now, observe that conditioned on $\mathcal F_k = \sigma(x_{k-1}, S_{k-1}, Z_{k-1})$, all the terms in the right hand side, except $e^{\lambda_{k+1}(M_{11}+M_{13})}$, are measurable. Thus, taking the conditional expectation in the above inequality, we get
    \begin{align*}
        &\E{e^{\lambda_{k+1}\lrp{M(x_{k+1} - x^*) + \alpha_k d_{k+1}}} \big| \mathcal F_k } \\
        &\,\,\,\le e^{\lambda_{k+1}M(\xk - x^*) \lrp{1 - 2\alpha_k\tilde{\gamma}_c + \alpha^2_k C_1 }} \E{e^{\lambda_{k+1} (M_{11}+M_{13})} \big| \mathcal F_k} e^{\lambda_{k+1}\alpha_{k-1}d_k (1-\alpha_k \tilde{\gamma}_c) + \lambda_{k+1}\alpha^2_k C_2} \\
        &\,\,\, \le e^{\lambda_{k+1}M(\xk - x^*) \lrp{1 - 2\alpha_k\tilde{\gamma}_c + \alpha^2_k C_1 }} \E{e^{\lambda_{k+1} M_{11}} \big| \mathcal F_k} \E{e^{\lambda_{k+1} M_{13}} \big| \mathcal F_k} e^{\lambda_{k+1}\alpha_{k-1}d_k (1-\alpha_k \tilde{\gamma}_c) + \lambda_{k+1}\alpha^2_k C_2},\numberthis\label{eq:boundcondexp}
    \end{align*}
    where the last inequality follows from the conditional independence of $M_{11}$ and $M_{13}$ (Lemma~\ref{lem:T11Properties}). Next, using Lemma~\ref{lem:T11Properties} and Hoeffding's inequality, we get
    \[ \E{ e^{\lambda_{k+1} M_{13}} \big| \mathcal F_k} \le e^{\lambda^2_{k+1} \frac{\alpha^2_k L^2_s B^2_2}{2 \lcs^4} \|x_k - x^*\|^2_c  } \le  e^{\lambda^2_{k+1} \frac{\alpha^2_k L^2_s B^2_2 u}{2 \lcs^4} M(x_k - x^*)  }. \]
    Again, from Lemma~\ref{lem:T11Properties}, we see that the conditional mean of $M_{11}$, conditioned on $\mathcal F_k$, is $0$. Moreover, since $\|x_k-x^*\|_c^2 \leq T_k$ almost surely, $M_{11}$ is bounded as below:
    \begin{align*} |M_{11}| &\le \frac{2\alpha_k L_s L_1}{\lcs^2}\|\xk- x^*\|^2_c + \frac{2\alpha_k L_s L_2}{\lcs^2}\|\xk-x^*\|_c \\
    &\le \frac{2\alpha_k L_s L_1}{\lcs^2}\|\xk- x^*\|_c T^\frac{1}{2}_k + \frac{2\alpha_k L_s L_2}{\lcs^2}\|\xk-x^*\|_c \\
    &\le \frac{2\alpha_k L_s}{\lcs^2}\|\xk- x^*\|_c  \lrp{L_1 T^\frac{1}{2}_k + L_2}. 
    \end{align*}
    Hence, using Hoeffding's lemma, we have
    \begin{align*}
    \E{e^{\lambda_{k+1} M_{11}} \big| \mathcal F_k} 
        &\le e^{\lambda^2_{k+1}\alpha^2_k \frac{2 L^2_s}{\lcs^4}\|x_k - x^*\|^2_c\lrp{L_1 T^\frac{1}{2}_k + L_2}^2}\\
        &\le e^{ \lambda^2_{k+1} \alpha^2_k 
        \frac{2L^2_s u}{\lcs^4} M(x_k - x^*) \lrp{L_1 T^\frac{1}{2}_k + L_2}^2 }\\
        &\le e^{ \lambda^2_{k+1} \alpha^2_k \frac{4 L^2_s u }{\lcs^4}   M(x_k - x^*) \lrp{L^2_1 T_k + L^2_2} },
    \end{align*}
    where the last inequality follows from $(a+b)^2 \le 2a^2 + 2b^2$, for all $a,b\in \R$.  Using the above bounds on conditional expectations of $M_{11}$ and $M_{13}$  in Equation~\eqref{eq:boundcondexp}, we have    
    \begin{align*}
        &\E{ e^{\lambda_{k+1}\lrp{M(x_{k+1} - x^*) + \alpha_k d_{k+1}}} \big| \mathcal F_k } \\
        &\qquad\le e^{\lambda_{k+1} M(\xk - x^*) \lrp{1 - 2\alpha_k\tilde{\gamma}_c + \alpha^2_k C_1 }} e^{\lambda^2_{k+1} \frac{\alpha^2_k L^2_s B^2_2 u}{2\lcs^4} M(\xk - x^*) } \\
        &\qquad\qquad e^{\lambda^2_{k+1} \alpha^2_k \frac{4 L^2_s u }{\lcs^4}   M(x_k - x^*) \lrp{L^2_1 T_k+ L^2_2}  + \lambda_{k+1}\alpha_{k-1}d_k (1-\alpha_k \tilde{\gamma}_c) + \lambda_{k+1}\alpha^2_k C_2}\\
        &\qquad \le  e^{\lambda_{k+1} M(\xk - x^*) \lrp{1 - 2\alpha_k\tilde{\gamma}_c + \alpha^2_k \lrp{  C_1 + \lambda_{k+1}  \frac{4 L^2_s u}{\lcs^4}\lrp{L^2_1 T_k + L^2_2 + B^2_2 } }}} \\
        &\qquad\qquad e^{\lambda_{k+1}\alpha_{k-1}d_k (1-\alpha_k \tilde{\gamma}_c) + \lambda_{k+1}\alpha^2_k C_2}\\
        &\qquad =  e^{ \lambda_{k+1}\lrp{M(\xk - x^*) M_3 + \alpha_{k-1}d_k (1-\alpha_k \tilde{\gamma}_c) }} e^{\lambda_{k+1} \alpha^2_k C_2}, \numberthis \label{eq:condbound1}
        \end{align*}
        where, 
        \[ M_3 := 1 - 2\alpha_k\tilde{\gamma}_c + \alpha^2_k \lrp{  C_1 + \lambda_{k+1}  \frac{ 4 L^2_s u}{\lcs^4}\lrp{L^2_1 T_k + L^2_2 + B^2_2 }  }. \]
    Now, from the definition of $\lambda_k$, we have
    \begin{align*}
        \alpha^2_k \lambda_{k+1}\frac{4 L^2_s u}{\lcs^4}\lrp{L^2_1 T_k + L^2_2 + B^2_2}
        &=    \frac{\theta 4 L^2_s u \alpha^2_k}{\alpha_{k+1}\lcs^4 T_{k+1}}\lrp{L^2_1 T_k + L^2_2 + B^2_2} \\
        &\le \frac{\theta 4 L^2_s u \alpha^2_k}{\alpha_{k+1}\lcs^4 T_{k+1}}\lrp{L^2_1 T_{k+1} + L^2_2 + B^2_2} \tag{$\because T_k\le T_{k+1}$}\\
        &\le \frac{\theta 4 L^2_s u \alpha^2_k}{\alpha_{k+1}\lcs^4 T_0}\lrp{L^2_1 T_0 + L^2_2 + B^2_2}\\
        &\le \frac{\theta 4 L^2_s u \alpha^2_k}{\alpha_{k+1}\lcs^4 \|x_0 - x^*\|^2_c}\lrp{L^2_1 \|x_0 - x^*\|^2_c + L^2_2 + B^2_2} \tag{$\because T_0 \ge \|x_0 - x^*\|^2_c$} \\
        &= \frac{\eta \alpha^2_k}{8 \alpha_{k+1}},
    \end{align*}
    where the last equality follows from choosing 
    \[ \theta =  \frac{\eta \lcs^4 \|x_0 - x^*\|^2_c}{32 L^2_s u  \lrp{L^2_1\|x_0 - x^*\|^2_c + L^2_2 + B^2_2} }. \]
    Next, since for $k \ge 0$ and $h \ge 8$, from Lemma~\ref{lem:ratioofalpha}, we have 
    \[ \frac{\alpha_{k}}{\alpha_{k+1}} \le 2, \]
    and hence, we further have
    \[ \alpha^2_k \lambda_{k+1}  \frac{4 L^2_s u}{\lcs^4}\lrp{L^2_1 T_k + L^2_2 + B^2_2} \le \frac{\eta \alpha_k}{4}, \]
    which implies 
    \[ M_3 \le  1 - \alpha_k \eta + \alpha^2_k C_1 + \frac{\eta\alpha_k}{4}. \]
    Now, choosing $h \ge 8$ large enough so that \[ \alpha_0  \le \frac{\eta}{4 C_1}, \]
    ensures that 
    \[ M_3 \le 1 - \alpha_k \eta + \frac{\alpha_k \eta}{4} + \frac{\alpha_k \eta}{4} = 1 - \frac{\alpha_k \eta}{2} = 1 - \alpha_k \tilde{\gamma}_c. \]

    Substituting this back in Equation~\eqref{eq:condbound1}, since $M(\cdot) \ge 0$, we have
    \begin{align*}
        &\E{ e^{\lambda_{k+1}\lrp{M(x_{k+1} - x^*) + \alpha_k d_{k+1}}} \big| \mathcal F_k }\\
        &\qquad \le e^{ \lambda_{k+1}(1-\alpha_k \tilde{\gamma}_c)\lrp{M(\xk - x^*) + \alpha_{k-1}d_k  }} e^{\lambda_{k+1} \alpha^2_k C_2}\\
        &\qquad \le e^{ \lambda_{k+1}(1-\alpha_k \tilde{\gamma}_c)\lrp{M(\xk - x^*) + \alpha_{k-1}d_k  }}  e^{2 \lambda_k \alpha^2_k C_2}, \numberthis\label{eq:expbound2}
    \end{align*}
    where the last inequality follows since $\alpha^2_k C_2 \ge 0$, and 
    \begin{equation*}
        \lambda_{k+1} = \lambda_k\frac{\lambda_{k+1}}{\lambda_k} = \lambda_k\frac{\alpha_k T_k}{\alpha_{k+1}T_{k+1}} \le \lambda_k\frac{\alpha_k}{\alpha_{k+1}} \le 2\lambda_k \tag{$\because   T_k \le T_{k+1}$ and Lemma~\ref{lem:ratioofalpha}}. 
    \end{equation*}
    Now, multiplying both sides of Equation~\eqref{eq:expbound2} by $e^{\lambda_{k+1}\alpha_k\frac{ L_s L_2}{4} }$, we get
\begin{align*}
        &\E{ e^{\lambda_{k+1}\lrp{M(x_{k+1} - x^*) + \alpha_k d_{k+1} + \alpha_k\frac{ L_s L_2}{4} }} \big| \mathcal F_k }\\
        &\qquad \le e^{ \lambda_{k+1}(1-\alpha_k \tilde{\gamma}_c)\lrp{M(\xk - x^*) + \alpha_{k-1}d_k + \alpha_{k-1}\frac{ L_s L_2}{4} } }  e^{2 \lambda_k \alpha^2_k C_2 + \lambda_{k+1} M_4}, \numberthis\label{eq:expbound3}
\end{align*}    
where,
\[ M_4 :=  \alpha_k\frac{ L_s L_2}{4}- (1-\alpha_k \tilde{\gamma}_c) \alpha_{k-1}\frac{ L_s L_2}{4}.\]
We now bound $M_4$ below. 
\begin{align*}
    M_4 &=  \frac{ L_s L_2}{4}\lrp{\alpha_k - \alpha_{k-1} + \alpha_k \alpha_{k-1}\tilde{\gamma}_c}\\
    &\le  \frac{ \tilde{\gamma}_cL_s L_2}{4}\alpha_k \alpha_{k-1} \tag{ $\because \alpha_k \le \alpha_{k-1}$}\\
    &= \alpha^2_k\frac{\alpha_{k-1}}{\alpha_k} \frac{ \tilde{\gamma}_cL_s L_2}{4}\\
    &\le 2 \alpha^2_k\frac{ \tilde{\gamma}_cL_s L_2}{4} \tag{Lemma~\ref{lem:ratioofalpha}}.
\end{align*}
Substituting the above in Equation~\eqref{eq:expbound3}, and as earlier, using $\lambda_{k+1} \le 2 \lambda_k$, we get
\begin{align*}
        &\E{ e^{\lambda_{k+1}\lrp{M(x_{k+1} - x^*) + \alpha_k d_{k+1} + \alpha_k \frac{L_s L_2}{4} }} \big| \mathcal F_k }\\
        &\qquad \le e^{ \lambda_{k+1}(1-\alpha_k \tilde{\gamma}_c)\lrp{M(\xk - x^*) + \alpha_{k-1}d_k + \frac{\alpha_{k-1} L_sL_2}{4} } }  e^{2 \lambda_k \alpha^2_k \lrp{C_2 + \frac{\tilde{\gamma}_c L_sL_2}{2} }}\\
        &\qquad \le \lrp{e^{ \lambda_{k}\lrp{ M(\xk - x^*) + \alpha_{k-1}d_k + \frac{\alpha_{k-1} L_sL_2}{4} }}}^{\frac{\lambda_{k+1}}{\lambda_k}(1-\alpha_k \tilde{\gamma}_c)}  e^{2 \lambda_k \alpha^2_k \lrp{C_2 + \frac{\tilde{\gamma}_c L_sL_2}{2} }}.
\end{align*}    

Next, from Lemma~\ref{lem:nonnegLyapunov}, we find that the exponent of $e$ is nonnegative in the above expression. Using $\lambda_{k+1}/\lambda_k \le \alpha_k/\alpha_{k+1}$ in the above expression, we get 
\begin{align*}
        &\E{ e^{\lambda_{k+1}\lrp{M(x_{k+1} - x^*) + \alpha_k d_{k+1} + \frac{\alpha_k L_sL_2 }{4 } }} \big| \mathcal F_k }\\
        &\qquad \le \lrp{e^{ \lambda_{k}\lrp{ M(\xk - x^*) + \alpha_{k-1}d_k + \frac{\alpha_{k-1} L_sL_2}{4}}}}^{\frac{\alpha_{k}}{\alpha_{k+1}}(1-\alpha_k \tilde{\gamma}_c)}  e^{2 \lambda_k \alpha^2_k \lrp{C_2 + \frac{\tilde{\gamma}_c L_sL_2}{2} }},%\numberthis\label{eq:expbound4}
\end{align*}
proving the first inequality.

Next, we prove that the exponent of the first term lies in $[0,1]$. Clearly, $(1-\alpha_k \tilde{\gamma}_c) \ge 0$ follows from choosing $h \ge 8$ large enough so that $\alpha_0 \le 1/\tilde{\gamma}_c$. Next, for proving the upper bound, consider the following for $z \in [0,1]$. 
\begin{align*}
    \frac{\alpha_k }{\alpha_{k+1}}(1-\alpha_k\tilde{\gamma}_c) 
    &= \lrp{\frac{k+h +1 }{k+h}}^z\lrp{1-\frac{\alpha \tilde{\gamma}_c}{(k+h)^z}}\\
    &\le \lrp{\frac{k+h+1}{k+h}}^ze^{-\frac{\alpha \tilde{\gamma}_c}{(k+h)^z}}\\
    &= \lrp{\lrp{1+\frac{1}{k+h}}^{k+h}}^\frac{z}{k+h} e^{-\frac{\alpha \tilde{\gamma}_c}{(k+h)^z}}\\
    &\le e^{ \frac{z}{k+h} - \frac{ \alpha \tilde{\gamma}_c}{(k+h)^z}} \tag{$\lrp{1+1/x}^x \le e$, for all $x \in \R$}.
\end{align*}

When $z=1$, on choosing $\alpha > 1/\tilde{\gamma}_c$, we have a negative exponent in the above expression, ensuring $(\alpha_{k}/\alpha_{k+1}) (1-\alpha_k \tilde{\gamma}_c) < 1$. 

When $z \in (0,1)$, choosing $h \ge \lrs{2z/(\alpha\tilde{\gamma}_c)}^{1/(1-z)}$, we have
\begin{align*}
    \frac{\alpha_k}{\alpha_{k+1}} \lrp{1 - \alpha_k \tilde{\gamma}_c} &\le e^{ \frac{z}{k+h} - \frac{ \alpha \tilde{\gamma}_c}{(k+h)^z}} \\
    &= e^{ \frac{z - (k+h)^{1-z}\alpha \tilde{\gamma}_c}{k+h}} \\
    &\le e^{ \frac{z - h^{1-z}\alpha \tilde{\gamma}_c}{k+h}}\\
    &\le 1.
\end{align*}
\end{proof}

\end{document}
